% Abbreviations and List of Symbols for the thesis
% Paste these sections where the corresponding sections appear in the main file.
% Required packages already present in the thesis preamble: longtable, array, amsmath.

% -------------------- Abbreviations --------------------
\section*{{\huge Abbreviations}\hfill}
\addcontentsline{toc}{chapter}{\textbf{Abbreviations}}
\vspace{0.5cm}

\begin{longtable}{p{3.2cm} p{11cm}}
6G & Sixth-generation mobile communication network \\
AI & Artificial intelligence \\
DDPG & Deep deterministic policy gradient \\
DL & Downlink \\
DT & Digital Twin \\
eMBB & Enhanced mobile broadband \\
GNN & Graph neural network \\
gNB & Next-generation Node B; a 5G/6G base station \\
GAE & Generalised advantage estimation \\
GRU & Gated recurrent unit \\
IPPO & Independent proximal policy optimisation \\
LLM & Large language model \\
MAE & Mean absolute error \\
MADDPG & Multi-agent deep deterministic policy gradient \\
MAPE & Mean absolute percentage error \\
MAPPO & Multi-agent proximal policy optimisation \\
MARL & Multi-agent reinforcement learning \\
MCP & Model Context Protocol \\
MLP & Multilayer perceptron \\
mMTC & Massive machine-type communication \\
MSE & Mean squared error \\
NDT & Network Digital Twin \\
NR & New Radio \\
O-RAN & Open Radio Access Network \\
PPO & Proximal policy optimisation \\
POMDP & Partially observable Markov decision process \\
PRB & Physical resource block \\
QoS & Quality of service \\
RAN & Radio access network \\
RL & Reinforcement learning \\
RMSE & Root mean squared error \\
SCARF & Self-supervised contrastive learning using random feature corruption \\
SLA & Service level agreement \\
SINR & Signal-to-interference-plus-noise ratio \\
SSL & Self-supervised learning \\
TGNN & Temporal graph neural network \\
UE & User equipment \\
URLLC & Ultra-reliable low-latency communication \\
VNF & Virtual network function \\
\end{longtable}

\newpage

% -------------------- List of Symbols --------------------
\section*{{\huge List of Symbols}\hfill}
\addcontentsline{toc}{chapter}{\textbf{List of Symbols}}
\vspace{0.5cm}

\begin{longtable}{p{3.2cm} p{11cm}}
$i$ & Index of a network cell or cell-level agent \\
$t$ & Time-step or traffic interval index \\
$\mathcal{V}$ & Set of cells (nodes) in the network graph \\
$\mathcal{E}$ & Set of edges representing connections or neighbourhood relations between cells \\
$\mathcal{G}=(\mathcal{V},\mathcal{E})$ & Network graph, consisting of cells and their relationships \\
$s_t^i$ & State or observation available to the agent controlling cell $i$ at time $t$ \\
$a_t^i$ & Allocation action selected by the agent for cell $i$ at time $t$ \\
$R_t^i$ & Reward received by the agent controlling cell $i$ at time $t$ \\
$S_t^i$ & Weighted service-level-agreement (SLA) score for cell $i$ at time $t$ \\
$F_t^i$ & Fraction of demand served by cell $i$ at time $t$ \\
$H_t^i$ & Reconfiguration churn for cell $i$ at time $t$ \\
$D_{t,k}^i$ & Demand for slice $k$ in cell $i$ at time $t$ \\
$\widehat{D}_{t+1,k}^i$ & Forecast demand for slice $k$ in cell $i$ for the next interval \\
$C_t^i$ & Usable capacity of cell $i$ at time $t$ \\
$x_{t,k}^i$ & Fraction of cell $i$'s capacity allocated to slice $k$ at time $t$ \\
$x_{t,\mathrm{idle}}^i$ & Fraction of cell $i$'s capacity left idle at time $t$ \\
$k$ & Index identifying a network slice, such as eMBB, URLLC or mMTC \\
$L$ & Network latency or packet/frame delay, depending on the reported experiment \\
$J$ & Jitter, representing variation in packet/frame delay \\
$\rho$ & Reliability or probability of successful service, when used in the SLA definition \\
$\mathbf{z}_t^i$ & Learned feature representation (embedding) for cell $i$ at time $t$ \\
$\widehat{\mathbf{y}}_t$ & Model-predicted value or traffic demand; the exact target depends on the experiment \\
\end{longtable}

% Notes:
% 1. The reward equation in Chapter 4 defines
%    R_t^i = 0.7 S_t^i + 0.3 F_t^i - 0.1 H_t^i.
% 2. Keep only symbols that are used in the final thesis text/equations.
% 3. If the thesis uses a different notation for any quantity, update the entry here
%    to match the notation used in the relevant chapter.
